% ============================================================
% CAPITOLO 3 -- TEST MANUALI (CATEGORY PARTITION)
% ============================================================
\section{Test manuali (Category Partition)}

La tecnica Category-Partition è un metodo di progettazione di test
black-box che procede in quattro fasi: identificazione dei domini di
input della funzionalità sotto test, a partire dalla documentazione o
dal comportamento atteso; derivazione, per ciascun dominio, delle
classi di equivalenza, ovvero sottoinsiemi di valori per cui ci si
aspetta un comportamento omogeneo; combinazione delle classi di
equivalenza individuate per i diversi parametri, per ottenere le
combinazioni di input da testare; eliminazione delle combinazioni
prive di senso o non raggiungibili nella pratica. Per ciascuna
combinazione ammissibile risultante viene infine progettato un singolo
caso di test.

Questa tecnica è stata applicata a entrambe le classi selezionate,
individuando per ciascuna un insieme di categorie corrispondenti alle
diverse configurazioni di comportamento osservabile. Le sezioni
seguenti descrivono, per ciascuna classe, la struttura interna
rilevante ai fini del testing, le categorie individuate e i risultati
ottenuti.

\subsection{La classe \texttt{Resolver}}
La classe \texttt{Resolver} è una classe che si occupa di determinare
come adattare i dati scritti in uno schema (\emph{writer}) a uno schema
di lettura (\emph{reader}) potenzialmente diverso. Implementa un
algoritmo in grado di confrontare ricorsivamente gli alberi degli schemi
e produrre un albero di azioni necessarie per la conversione dei dati.

Il risultato prodotto è un'istanza della classe astratta
\texttt{Resolver.Action}, il cui tipo concreto a runtime determina la
specifica trasformazione da applicare: \texttt{Action} non viene mai
istanziata direttamente, ma sempre tramite una delle sue sottoclassi.
La \tabref{tab:resolver-action-categorie} riassume quali sono i
possibili tipi di azioni che la classe può restituire.

\subsection{Test manuali su \texttt{Resolver}}

Il metodo \texttt{resolve()} accetta due parametri principali:
lo schema del \emph{writer} e lo schema del \emph{reader} (un terzo
parametro opzionale, \texttt{GenericData}, interviene solo nel recupero
di valori di default e conversioni logiche, e non introduce ulteriori
categorie di comportamento rilevanti ai fini del testing). Il criterio
di testing adottato si è basato sull'interazione tra le diverse
combinazioni possibili di questi due parametri. Sono state individuate
otto categorie di test, ciascuna corrispondente a una delle otto
sottoclassi di \texttt{Resolver.Action} discusse in precedenza.

L'approccio adottato è prevalentemente black-box, dato che le regole di
comportamento attese (quali promozioni numeriche sono ammesse, come si
risolvono record, enum e union) derivano dalla specifica pubblica di
schema evolution di Avro. Tuttavia, la \emph{verifica} del risultato in
ciascun test avviene tramite l'operatore \texttt{instanceof} su nomi di
classi concrete interne (ad esempio \texttt{Resolver.DoNothing},
\texttt{Resolver.Promote}), informazione non derivabile dalla sola
documentazione esterna: la progettazione dei casi di test è quindi
black-box, mentre la loro implementazione richiede un elemento di
conoscenza white-box del codice sorgente.

\begin{itemize}
    \item \textbf{DoNothing}: verificata su tutti gli otto tipi
    primitivi di Avro, con schemi identici tra writer e reader.
    \item \textbf{Promote}: verificato che le promozioni numeriche
    valide (ad esempio \texttt{INT}$\to$\texttt{LONG},
    \texttt{FLOAT}$\to$\texttt{DOUBLE}) producano l'azione corretta, e
    che le promozioni non valide generino invece un errore.
    \item \textbf{ErrorAction}: verificato che le incompatibilità
    irrisolvibili tra i due schemi generino correttamente un errore.
    \item \textbf{Container}: verificato che, per \texttt{ARRAY} e
    \texttt{MAP}, la risoluzione ricorsiva degli elementi interni
    funzioni correttamente.
    \item \textbf{EnumAdjust}: verificato che la rimappatura dei
    simboli tra writer e reader funzioni correttamente.
    \item \textbf{RecordAdjust}: verificato che la risoluzione campo
    per campo del record funzioni correttamente.
    \item \textbf{ReaderUnion}: verificato che la ricerca del ramo
    della union compatibile con il tipo scritto funzioni correttamente.
    \item \textbf{WriterUnion}: verificato che la risoluzione di ogni
    possibile ramo scritto dal writer rispetto al reader funzioni
    correttamente.
\end{itemize}

Per alcune categorie, in particolare \texttt{Promote}, i valori
scelti per i test non sono stati selezionati arbitrariamente
all'interno della categoria, ma rappresentano i confini tra
comportamento valido e non valido: per ogni promozione numerica
valida testata (ad esempio \texttt{INT}$\to$\texttt{LONG}), è stata
verificata anche la direzione opposta, non valida
(\texttt{LONG}$\to$\texttt{INT}), proprio in corrispondenza del limite
tra le due categorie. La \tabref{tab:resolver-boundary-values} riporta,
per ciascuna categoria, i valori concreti utilizzati nella
progettazione dei test.

In totale sono stati scritti ed eseguiti 32 test manuali, tutti
passanti: non è emersa alcuna anomalia comportamentale in questa fase.
La \tabref{tab:resolver-elenco-test} riporta l'elenco completo dei 32
test, con l'input e l'assert atteso per ciascuno. I test sono
collocati nel file \texttt{TestResolverCategoryPartition.java},
all'interno del package \texttt{org.apache.avro} del progetto.


\subsection{Copertura strutturale di \texttt{resolver}}

I dati di copertura sono stati misurati con JaCoCo, isolando
l'esecuzione della sola suite di test manuali. La
\tabref{tab:bb-coverage} riporta i valori di \emph{statement coverage}
(42.6\%) e \emph{branch coverage} (31.6\%) ottenuti; la
\tabref{fig:jacoco-resolver-bb} mostra il dettaglio per metodo.

La copertura risultante è parziale: il category-partition, per sua
natura, individua le categorie a partire dalle funzionalità e dal
comportamento atteso della classe, non dall'obiettivo di massimizzare
la copertura di ogni singola riga o diramazione del codice. Una
copertura non completa non indica quindi una debolezza della suite, ma
una caratteristica intrinseca della tecnica adottata: la relazione tra
copertura strutturale e adeguatezza dei test manuali rispetto alle
altre tecniche impiegate viene approfondita nel
Capitolo~\ref{sec:adeguatezza}, mentre i mutanti sopravvissuti nelle
porzioni di codice non coperte sono discussi nel
Capitolo~\ref{sec:mutation}.

\subsection{La classe \texttt{Utf8}}

La classe \texttt{Utf8} si occupa di rappresentare una sequenza di
testo direttamente come byte codificati in UTF-8. In Java, il tipo
\texttt{String} è rappresentato internamente come una sequenza di
caratteri UTF-16: se un dato testuale dovesse essere sempre convertito
da \texttt{String} a byte UTF-8 e viceversa, si introdurrebbe un costo
computazionale ripetuto e spesso non necessario. La classe \texttt{Utf8}
risolve questo problema rappresentando direttamente il testo come byte
UTF-8, evitando conversioni superflue e convertendo in \texttt{String}
solo quando effettivamente richiesto.

La classe implementa tre interfacce:
\begin{itemize}
    \item \texttt{CharSequence}: consente di trattare l'istanza come
    se fosse testo, fornendo metodi quali \texttt{length()},
    \texttt{charAt()} e \texttt{subSequence()};
    \item \texttt{Comparable<Utf8>}: permette il confronto tra due
    istanze ai fini dell'ordinamento;
    \item \texttt{Externalizable}: fornisce un meccanismo di
    serializzazione a basso livello, tramite i metodi
    \texttt{writeExternal} e \texttt{readExternal}.
\end{itemize}

Gli attributi principali della classe sono: un array di byte
(\texttt{bytes}), contenente la rappresentazione UTF-8 del testo; un
intero (\texttt{length}), che indica quanti byte dell'array sono
effettivamente utilizzati (l'array può essere più grande della
lunghezza logica del contenuto, per consentire il riutilizzo della
stessa area di memoria a fronte di modifiche successive); una cache
opzionale della \texttt{String} corrispondente; una cache opzionale
dell'hash calcolato.

La \tabref{tab:utf8-metodi} elenca i metodi pubblici principali della
classe.

\subsection{Test manuali su \texttt{Utf8}}

La classe \texttt{Utf8} è composta da più metodi indipendenti: sono
state definite sei categorie di test, raggruppate per area funzionale
correlata.

\begin{itemize}
    \item \textbf{Costruzione}: verificata la corretta inizializzazione
    dell'istanza a partire da ciascuna delle modalità di creazione
    disponibili (costruttore vuoto, da stringa vuota, ASCII e Unicode,
    da stringa \texttt{null}, di copia, da array di byte vuoto e
    ASCII).
    \item \textbf{Uguaglianza e ordinamento}: verificato il
    comportamento di \texttt{equals()} e \texttt{compareTo()} rispetto
    alla stessa istanza, a contenuto identico, a contenuto diverso a
    parità di lunghezza, a lunghezza diversa, e nel confronto con
    \texttt{null} o con un oggetto di tipo diverso. Per questa
    categoria, l'approccio adottato non è stato puramente black-box:
    \texttt{equals()} e \texttt{hashCode()} implementano internamente
    due percorsi di calcolo distinti, uno ottimizzato per array
    esattamente dimensionati e uno tramite loop manuale, con soglia
    di attivazione a 7 byte. L'individuazione di questo confine ha
    richiesto un'ispezione mirata del codice sorgente, introducendo
    un elemento di conoscenza white-box a supporto della
    progettazione dei casi di test limite.
    \item \textbf{Gestione della lunghezza}: verificato il
    comportamento di \texttt{setByteLength()} in aumento, in
    diminuzione, a lunghezza invariata e nel caso limite di
    azzeramento.
    \item \textbf{\texttt{CharSequence}}: verificato il comportamento
    di \texttt{charAt()} (primo carattere, ultimo carattere, indice
    fuori intervallo) e di \texttt{subSequence()} (intervallo vuoto,
    valido, non valido).
    \item \textbf{Serializzazione}: verificato che una coppia di
    operazioni \texttt{writeExternal}/\texttt{readExternal}
    (round-trip) preservi correttamente il contenuto, su un'istanza
    vuota, su contenuto ASCII e su contenuto Unicode.
    \item \textbf{\texttt{getBytesFor}}: verificato il comportamento
    del metodo statico di utilità su stringa vuota, stringa ASCII,
    stringa Unicode e argomento \texttt{null}.
\end{itemize}

La \tabref{tab:utf8-elenco-test} riporta l'elenco completo dei 35 test,
con l'input/azione e l'assert atteso per ciascuno.

In totale sono stati scritti ed eseguiti 35 test manuali. La scoperta
emersa dal test relativo alla gestione della lunghezza in aumento è
discussa nel dettaglio nella sezione successiva.
\subsection{Scoperta: anomalia di \texttt{setByteLength}}

Il test T19, relativo alla gestione della lunghezza in aumento tramite
il metodo \texttt{setByteLength()}, ha evidenziato un'anomalia nel
comportamento della classe \texttt{Utf8}. In particolare, quando si
aumenta la lunghezza logica di un'istanza già popolata, ci si
aspetterebbe che il contenuto originale venga preservato e che i nuovi
byte aggiunti siano inizializzati a zero. Il test ha invece mostrato
che il contenuto originale viene azzerato, restituendo l'intero buffer
vuoto anziché esteso.

La scoperta è avvenuta empiricamente: l'assert inizialmente scritto,
basato sull'aspettativa di preservazione del contenuto, falliva
all'esecuzione. Il test è stato mantenuto, riformulato per verificare
il comportamento realmente osservato, poiché ha permesso di scoprire
questo comportamento anomalo, che sarebbe altrimenti passato
inosservato.



\subsection{Copertura strutturale di \texttt{utf8}}
I dati di copertura sono stati misurati con JaCoCo, isolando
l'esecuzione della sola suite di test manuali. La
\tabref{tab:bb-coverage} riporta i valori di \emph{statement coverage}
(66.7\%) e \emph{branch coverage} (43.2\%) ottenuti; la
\tabref{fig:jacoco-utf8-bb} mostra il dettaglio per metodo.

Anche in questo caso la copertura risulta parziale, per lo stesso
motivo discusso per \texttt{Resolver}: le categorie individuate
rappresentano il comportamento funzionale atteso della classe, non un
obiettivo di copertura strutturale completa.